\documentclass[11pt]{article}
\usepackage{amsmath,amssymb,amsthm,geometry,hyperref}
\geometry{margin=1in}
\newtheorem{theorem}{Theorem}
\title{A formal certificate for a balanced-matrix discrepancy counterexample}
\author{Myeongjin Kang}
\date{October 25, 2026}
\begin{document}
\maketitle
\begin{abstract}
We give a Lean certificate for the counterexample previously posted by Shivam
Patel to MathDB340783. Equal row and column energies do not force discrepancy
to propagate from one row to all rows. The certificate includes positivity,
every positive discrepancy tolerance, and every reflexive interpretation of
the source's approximate-energy comparison. No mathematical priority is claimed.
\end{abstract}
For a real $n\times m$ matrix $B$, put
\[
 R_i(B)=\sum_j b_{ij}^2,\qquad C_j(B)=\sum_i b_{ij}^2,
 \qquad M_i(B)=m^{-1}\sum_j b_{ij}.
\]
A row $i$ is $\epsilon$-fair if $|M_i(B)-b_{ij}|<\epsilon$ for every $j$.
Definition 3.1 of \cite{agama} compares the row energies to each other and
the column energies to each other using an unspecified approximate relation.
Conjecture 6.6 and Remark 6.7 propose propagation of fairness from one row
to all rows. We work with exact equality, which satisfies any reflexive
interpretation of that approximate comparison.
\begin{theorem}
For every $\epsilon>0$ there is a strictly positive nonzero real $3\times3$
matrix with all row and column energies equal, whose first row is
$\epsilon$-fair and whose second row is not. The same conclusion holds if
energy equality is replaced by any prescribed reflexive relation on real
numbers, or by any positive absolute energy tolerance.
\end{theorem}
\begin{proof}
Use $B=\epsilon A$, where the witness from \cite{patel} is
\[
A=\begin{pmatrix}5&5&5\\1&5&7\\7&5&1\end{pmatrix}.
\]
Every row and every column of $A$ has squared energy $75$, so those of
$B$ have energy $75\epsilon^2$. Positivity and nonzeroness follow from
$\epsilon>0$. The first row of $B$ is constant. The second row has mean
$13\epsilon/3$, and its first deviation is
\[
 |13\epsilon/3-\epsilon|=10\epsilon/3>\epsilon.
\]
Equal energies satisfy any reflexive relation. Their difference is zero,
so it is smaller in absolute value than any positive energy tolerance.
\end{proof}
\paragraph{Formal coverage.}
The accompanying \texttt{Discrepancy.lean} proves the theorem and the
explicit negation of the universal exact-balance propagation statement.
The final declaration \texttt{every\_reflexive\_interpretation\_counterexample}
exposes the only semantic requirement on the undefined approximate relation.
The build checks all local types and axioms using Lean 4.34.1 and mathlib
revision \texttt{d13f23b723b8a846827a245b89c10fc7d3f11612}.
Only standard kernel axioms occur. The compiler and precompiled dependency
objects are trusted. This is a complete certificate for this counterexample,
not an assertion that an unspecified nonreflexive relation was formalized.
\paragraph{Attribution.}
The matrix and scaling argument were already in Patel's MathDB reply dated
August 22, 2026. The contribution here is the formal certificate and the
explicit relation-parametric treatment. No submission or external review is
claimed. Original prose is CC BY 4.0; original Lean code is MIT.
\begin{thebibliography}{9}
\bibitem{agama} T. Agama and G. Kibiti,
\emph{Balanced matrices}, arXiv:1810.07542v3, Definition 3.1,
Conjecture 6.6, Remark 6.7.
\url{https://arxiv.org/html/1810.07542v3}.
\bibitem{patel} S. Patel, Counterexample reply to MathDB340783,
August 22, 2026, reply identifier
\texttt{be21f974-e069-4a4c-b446-252e228a1119}.
\url{https://mathdb.com/p/340783/the-fair-discrepancy-propagation-conjecture-for-fully-balanc}.
\end{thebibliography}
\end{document}
